\section{Enquadramento Teórico}
\label{sec:background}

Inclui fórmulas, definições e referências.

Exemplo de equação:
\begin{equation}
  E = mc^2
\end{equation}

Fórmula resolvente de uma equação do segundo grau $ax^2 + bx + c = 0$, com $a \neq 0$:
\begin{equation}
  x = \frac{-b \pm \sqrt{b^2 - 4ac}}{2a}
  \label{eq:resolvente}
\end{equation}
O sinal do discriminante $\Delta = b^2 - 4ac$ determina o número de soluções reais (ver Apêndice~\ref{apx:extra}).

Outros exemplos simples:
\begin{align}
  \sum_{i=1}^{n} i &= \frac{n(n+1)}{2} \\
  \int_0^1 x^2 \, \mathrm{d}x &= \frac{1}{3} \\
  \lim_{n \to \infty} \left(1 + \frac{1}{n}\right)^n &= e
\end{align}

Exemplo de tabela:
\begin{table}[H]
  \centering
  \caption{Exemplo de tabela}
  \begin{tabular}{@{}lll@{}}
    \toprule
    \textbf{Item} & \textbf{Valor} & \textbf{Notas} \\
    \midrule
    A & 1.0 & Texto \\
    B & 2.0 & Texto \\
    \bottomrule
  \end{tabular}
\end{table}
